\clearpage
\section{Suy diễn mẫu lớn và cách chọn phương pháp}
\label{sec:large}

\subsection{Điểm thay đổi: phân phối giới hạn thay cho phân phối chính xác}
Trong tổng thể không chuẩn, $T^2$ vẫn tính được nếu $S$ khả nghịch,
nhưng công thức $F$ của mẫu chuẩn không còn là một kết quả chính xác.
Với mẫu iid, $p$ cố định, trung bình và hiệp phương sai hữu hạn, $\Sigma$
xác định dương, Phần~\ref{sec:sampling} đã cho
\[
 n(\Xbar-\vmu)'S^{-1}(\Xbar-\vmu)\xrightarrow{d}\chi^2_p.
\]
Suy diễn mẫu lớn dùng phân phối giới hạn này để xấp xỉ ngưỡng.
\nguon{Chương 4, §4.5, (4-28), tr.~176; Chương 5, §5.5, (5-31), tr.~235}

\begin{result}[Kiểm định xấp xỉ]
Với các điều kiện trên và $n$ đủ lớn, bác bỏ $H_0:\vmu=\vmu_0$ ở mức
xấp xỉ $\alpha$ khi
\begin{equation}
 T^2=n(\xbar-\vmu_0)'S^{-1}(\xbar-\vmu_0)>\chi^2_{p;1-\alpha}.
 \label{eq:largetest}
\end{equation}
Giá trị p xấp xỉ là đuôi phải của $\chi^2_p$ tại $T^2$ quan sát.
\end{result}
\nguon{Chương 5, §5.5, Result 5.4, tr.~235}

Thống kê không đổi, nhưng ngưỡng thay đổi. Nếu tổng thể chuẩn và có đủ
điều kiện, có thể tiếp tục dùng ngưỡng $F$ chính xác dù mẫu lớn.
Không bắt buộc bỏ kết quả chính xác để dùng xấp xỉ.

\subsection{Khoảng tin cậy xấp xỉ}
Miền cho trung bình dùng
\[
 n(\xbar-\vmu)'S^{-1}(\xbar-\vmu)\le\chi^2_{p;1-\alpha}.
\]
Tương ứng, các khoảng đồng thời cho mọi tổ hợp là
\begin{equation}
 \mathbf a'\xbar\pm\sqrt{\chi^2_{p;1-\alpha}}
 \sqrt{\frac{\mathbf a'S\mathbf a}{n}}.
 \label{eq:largeinterval}
\end{equation}
Với từng thành phần, đặt $\mathbf a=\mathbf e_i$; với hiệu hai thành
phần, đặt $\mathbf a=\mathbf e_i-\mathbf e_k$ như trước.
\nguon{Chương 5, §5.5, Result 5.5, tr.~235--236}

Cho $m$ tổ hợp xác định trước, Bonferroni mẫu lớn dùng
\begin{equation}
 \mathbf a_j'\xbar\pm z_{1-\alpha/(2m)}
 \sqrt{\frac{\mathbf a_j'S\mathbf a_j}{n}},\qquad j=1,\ldots,m.
 \label{eq:largebonf}
\end{equation}
Đây là hệ quả của CLT một chiều cho từng tổ hợp và bất đẳng thức hợp.
$z_\gamma$ là phân vị của $\N(0,1)$.
Trong tổng thể không chuẩn, việc giữ hệ số $t$ không tự làm bảo đảm trở
thành chính xác; cả cách dùng $t$ lẫn $z$ đều cần đánh giá xấp xỉ.

\subsection{Một ví dụ để thấy độ chênh giữa hai ngưỡng}
\begin{example}[Số tóm tắt giả định, $n=100$]
Giả sử $p=2$, $n=100$, $\xbar=(6,10)'$ và
$S=\begin{pmatrix}8&-10/3\\-10/3&2\end{pmatrix}$.
Đây là một mẫu giả định mới; không lấy bốn quan sát rồi tự đổi $n$ thành 100.
Kiểm định $\vmu_0=(7,11)'$ cho
\[
 T^2=100\frac{75}{22}=340.9091.
\]
Ngưỡng chi-bình phương 95\% là 5.9915; nếu mẫu chuẩn, ngưỡng $F$ chính
xác là
\[
 \frac{2\cdot99}{98}F_{2,98;0.95}\approx6.2415.
\]
Cả hai cách đều bác bỏ ở ví dụ này.
Các khoảng thành phần chi-bình phương là
\[
 \mu_1:6\pm0.6923,\qquad\mu_2:10\pm0.3462.
\]
Với $m=3$ phát biểu định trước, hệ số Bonferroni chuẩn là
$z_{1-0.05/6}\approx2.3940$.
\end{example}

Hãy so với mẫu $n=4$ ở phần Hotelling: trung bình và $S$ có thể giống nhau,
nhưng sai số chuẩn của trung bình và ngưỡng suy diễn khác nhau. Vì vậy
không thể diễn giải độ chắc chắn mà quên kích thước mẫu.

\subsection{``Mẫu lớn'' không có một ngưỡng đúng cho mọi dữ liệu}
Một quy tắc thực hành thường gặp là $n-p\ge30$. Nó có thể hữu ích cho
ôn tập hoặc chỉ dẫn sơ bộ, nhưng không phải giả thiết của CLT hay Slutsky,
và không bảo đảm xấp xỉ tốt cho mọi phân phối.
Sách thảo luận rằng độ chính xác còn phụ thuộc số chiều và dạng phân phối;
hai bài toán cùng $n-p$ vẫn có thể khác nhau nhiều.
\nguon{Chương 5, §5.5, thảo luận về kích thước mẫu và lệch chuẩn, tr.~236}

Các yếu tố cần xem gồm độ lệch, đuôi nặng, ngoại lệ, sự phụ thuộc giữa
quan sát, tỉ lệ giữa $p$ và $n$, và tình trạng gần suy biến của $S$.
Giới hạn với $p$ cố định không tự áp dụng khi $p$ tăng nhanh cùng $n$.
Nếu phương sai không hữu hạn hoặc dữ liệu phụ thuộc thời gian, các công
thức giới hạn trong bài giảng không tự được bảo đảm.

\luuy{Mẫu lớn giúp trung bình mẫu gần chuẩn; nó không sửa được sai
thiết kế lấy mẫu, không làm quan sát phụ thuộc trở thành iid, và không
biến ma trận gần suy biến thành một ước lượng tốt.}

\subsection{Kiểm tra mô hình trước khi suy diễn}
Đối với dữ liệu thực, nên kiểm tra:
\begin{enumerate}
 \item \textbf{Thiết kế mẫu:} các hàng có phải những đối tượng độc lập,
 cùng cơ chế lấy mẫu hay có đo lặp, chuỗi thời gian, phân cụm?
 \item \textbf{Chất lượng dữ liệu:} đơn vị, giá trị thiếu, cột hằng,
 biến trùng và phép biến đổi đã dùng.
 \item \textbf{Đồ thị:} histogram, Q--Q plot cho từng biến; scatter plot
 cho các cặp; điểm ngoại lệ và các dạng phi tuyến.
 \item \textbf{Hiệp phương sai:} $n>p$, hạng và mức độ gần suy biến.
\end{enumerate}
Các Q--Q plot biên nhìn đẹp không chứng minh chuẩn đồng thời, như phản
ví dụ ở Phần~\ref{sec:linear}. Đồ thị là công cụ đánh giá, không thay thế
toàn bộ giả thiết. Khi đã biến đổi dữ liệu, kết luận thường nói về trung
bình \emph{trên thang biến đổi}; ví dụ $\E(\log X)$ không bằng
$\log\E(X)$ nói chung.
\nguon{Chương 4, §4.6, tr.~177--186; §4.7, tr.~187--192; biến đổi dữ liệu ở §4.8, từ tr.~192}

Không tự xóa ngoại lệ chỉ để đạt kết quả mong muốn. Cần xác định đó là
lỗi nhập liệu hay quan sát có ý nghĩa, rồi nêu rõ xử lý. Nếu cần phương
pháp cho dữ liệu phụ thuộc, hiệp phương sai suy biến hoặc cao chiều, phải
xây dựng mô hình khác; chúng nằm ngoài phạm vi bài giảng này.

\subsection{Bảng chọn phương pháp}
\begin{center}
\small
\begin{tabular}{p{.30\linewidth}p{.30\linewidth}p{.31\linewidth}}
\toprule
Câu hỏi & Điều kiện chính & Công cụ\\
\midrule
Phân phối một tổ hợp hoặc vectơ con & Vectơ gốc chuẩn & $A\vmu$, $A\Sigma A'$; chọn khối\\
Phân phối khi biết một khối & Chuẩn chung, khối đã biết khả nghịch & Công thức có điều kiện\\
Độc lập giữa các tổ hợp & Chuẩn đồng thời & Kiểm tra $\mathbf a'\Sigma\mathbf b=0$\\
Kiểm định vectơ trung bình chính xác & Mẫu iid chuẩn, $n>p$, $\Sigma$ xác định dương & $T^2$ và ngưỡng $F$\\
Mọi tổ hợp của trung bình & Như trên & Khoảng đồng thời $T^2$\\
Ít tổ hợp định trước & Mẫu iid chuẩn & Bonferroni với $m$ phát biểu\\
Trung bình, tổng thể không chuẩn & Iid, moment hữu hạn, $p$ cố định, mẫu đủ lớn & Ngưỡng $\chi^2$; Bonferroni $z$\\
\bottomrule
\end{tabular}
\end{center}

\subsection{Liên hệ với BTVN 1--2}
\begin{center}
\small
\begin{tabular}{lp{.68\linewidth}}
\toprule
Bài & Phần cần đọc và thao tác chính\\
\midrule
2.1 & Phần 3: chọn vectơ trung bình và ma trận hiệp phương sai con.\\
2.2 & Phần 3: kiểm tra khối hiệp phương sai hoặc hiệp phương sai của tổ hợp.\\
2.3 & Phần 4: phân hoạch theo biến mục tiêu và biến đã biết.\\
2.4a,b & Phần 3: tổ hợp tuyến tính; phương trình triệt tiêu hiệp phương sai.\\
2.9a--d & Phần 5 và 8: phân phối trung bình mẫu, $T^2$ và xấp xỉ; chọn theo giả thiết của từng câu.\\
3.1 & Phần 6: tính $T^2$, đổi sang $F$, so ngưỡng và kết luận.\\
3.2 & Phần 6: biến đổi đồng bộ dữ liệu, $S$ và trung bình giả thuyết.\\
3.3--3.5 & Phần 7: bảo đảm đồng thời, miền elip, khoảng thành phần và các hiệu; đếm đúng $m$.\\
3.6a--d & Phần 8: suy diễn mẫu lớn, khoảng $\chi^2$ và Bonferroni; nhớ hạng hiệp phương sai.\\
\bottomrule
\end{tabular}
\end{center}
Số bài ở đây theo BTVN địa phương. Chúng không phải số Result hay số
Exercise trong giáo trình. Khi đối chiếu lời giải, phải dùng đúng bộ dữ
kiện của từng bài; không suy rằng cùng một ký hiệu $\Sigma$ nghĩa là cùng
một ma trận ở mọi trang.

\kiemtra{Nếu tổng thể chuẩn và $n=100$, có bắt buộc dùng ngưỡng
$\chi^2$ không? Không. Ngưỡng $F$ chính xác vẫn dùng được.}
